\clearpage
\VerdierPhysicalPageAnchor{0210}%
\begingroup\small\setlength{\parskip}{0pt}\setlength{\emergencystretch}{1.5em}
\setlength{\abovedisplayskip}{4pt plus 2pt minus 2pt}%
\setlength{\belowdisplayskip}{4pt plus 2pt minus 2pt}%

\VerdierSemanticAnchor{ch3d-3-1-4-proof}
Only the third assertion remains to be proved. The other assertions are
immediate consequences of Proposition (2.4.4). Assertion 2) shows at once
that the cohomological dimension of $\mathcal A$ is less than or equal to
the supremum of the injective (respectively, projective) dimensions of the
objects of $\mathcal A$. Suppose that the cohomological dimension of
$\mathcal A$ is equal to $n$, and let

\VerdierSemanticAnchor{ch3d-3-1-4-proof-resolution}
\[
I^\bullet=\cdots\longrightarrow0\longrightarrow0\longrightarrow I^0
\longrightarrow I^1\longrightarrow I^2\longrightarrow I^3
\longrightarrow\cdots
\]
be a resolution of injective type of an object $Y$ of $\mathcal A$. For every
object $X$ of $\mathcal A$, we have
$\operatorname{Hom}_{K(\mathcal A)}(X[-n-1],I^\bullet)=0$. In particular,
taking $X=\operatorname{Ker}d_{I^\bullet}^{n+1}$, we see at once that the
canonical morphism
$\operatorname{Ker}d_{I^\bullet}^{n+1}\longrightarrow I^{n+1}$ factors
through $I^n$:

\VerdierSemanticAnchor{ch3d-3-1-4-proof-factorization}
\[
\begin{tikzcd}[column sep=4.8em,row sep=2.8em]
{} & \operatorname{Ker}d_{I^\bullet}^{n+1}
  \arrow[dl,dashed] \arrow[d] & {} \\
I^n \arrow[r] & I^{n+1} \arrow[r] & \cdots
\end{tikzcd}
\]

\VerdierSemanticAnchor{ch3d-3-1-4-proof-truncation}
Consequently, $\operatorname{Ker}d_{I^\bullet}^{n+1}$ and
$\operatorname{Ker}d_{I^\bullet}^{n}$ are direct summands of $I^n$ and are
therefore injective. The complex
\[
\cdots\longrightarrow0\longrightarrow I^0\longrightarrow I^1
\longrightarrow\cdots\longrightarrow I^{n-1}
\longrightarrow\operatorname{Ker}d_{I^\bullet}^{n}
\longrightarrow0\longrightarrow\cdots
\]
is an injective resolution of $Y$; this shows that the injective dimension of
$Y$ is less than or equal to $n$. It is proved similarly that the projective
dimension of every object $X$ of $\mathcal A$ is less than or equal to $n$.

\VerdierSemanticAnchor{ch3d-3-1-5}
\medskip\noindent\textbf{3.1.5.}\quad
Let $X$ and $Y$ be two objects of $\mathcal A$. Proposition (3.1.4) shows
that when the category $\mathcal A$ has enough injective or projective
objects, the groups $\operatorname{Ext}_{\mathcal A}^{i}(X,Y)$ defined here
coincide with those defined classically.

\VerdierSemanticAnchor{ch3d-3-2}
\subsection{The Yoneda construction.}

\VerdierSemanticAnchor{ch3d-3-2-1}
\noindent\textbf{3.2.1.}\quad
Let $X$ and $Y$ be two objects of $\mathcal A$. Denote by
$E_{\mathcal A}^{n}(X,Y)$ ($n\geq1$) the set of exact sequences

\VerdierSemanticAnchor{ch3d-3-2-1-sequence}
\[
S^n(X,Y)=0\longrightarrow Y\longrightarrow Z_{n-1}
\longrightarrow Z_{n-2}\longrightarrow\cdots\longrightarrow Z_0
\xrightarrow{u_0}X\longrightarrow0.
\]
The integer $n$ is called the length of the sequence $S^n(X,Y)$. Associate
with every element $S^n(X,Y)$ of $E_{\mathcal A}^{n}(X,Y)$ three complexes:
\endgroup
